\documentclass[sigconf,anonymous,nonacm]{acmart}
\settopmatter{printacmref=false,printccs=false,printfolios=true}
\renewcommand\footnotetextcopyrightpermission[1]{}
\citestyle{acmauthoryear}
\usepackage{amsmath,booktabs,tabularx}
\usepackage{siunitx,xurl,placeins}
\setlength{\emergencystretch}{2em}
\newtheorem{proposition}{Proposition}
\sisetup{detect-all,table-number-alignment=center}
\makeatletter\@ACM@balancefalse\makeatother
\raggedbottom
\acmConference[ICAIF 2026 Workshop]{The 3rd Workshop on LLMs and Generative AI for Finance}{November 2026}{Milan, Italy}
\title[Reasoning Budgets for LLM Decision Services]{Evaluating Reasoning Budgets for LLM Decision Services: Auction Utility, Calculator Gaps, and Resource Accounting}
\author{Anonymous Author(s)}
\begin{document}
\begin{abstract}
We present an evaluation design for numerical LLM decision services that connects held-out economic scoring, a construction-calculator comparator and recorded resource use. Six checkpoints choose bids in 216 constructed consumer-auction tasks under supplied win probabilities and assigned reasoning budgets. Increasing budgets from 0 to 768 tokens closes 42--85\% of the aggregate zero-to-calculator utility gap, while every checkpoint remains below the calculator in simultaneous dollar-profit comparisons. Observable action classes distinguish displayed-objective optimality from held-out performance. Ratio and dollar aggregation can disagree, gains are concentrated, and free-form validity varies across checkpoints; conditional-on-valid-output utility therefore requires explicit denominators. Stage-specific accounting separates generation/answer, readout and confidence-query resources while disclosing excluded overheads. An appendix analyzes an assumed service contest: pooled cost scenarios imply private payment break-evens of \$0.03--\$0.18 and planner coefficients of \$0.02--\$0.12, with hardware and product sensitivity. The contribution is a controlled auction benchmark and audit design, not a validated financial deployment or agent system.
\end{abstract}
\keywords{LLM decision services, numerical consistency, calculator comparator, held-out economic utility, test-time compute, resource accounting}
\maketitle
\section{Introduction}
An LLM embedded in a decision workflow should be assessed by the consequences of the action it recommends, as well as by the resources required to produce it. Fluent reasoning is insufficient when an inexpensive numerical method can solve the calculation. A useful preliminary evaluation isolates action selection, supplies a transparent objective and checks whether additional assigned reasoning budget buys economic utility.

A numerical subtask that can arise in financial decision-support services is choosing among candidate quotes or prices conditional on supplied probabilities. This is a motivating analogy: our benchmark isolates the calculation/action-selection step in a controlled consumer-auction setting. Models receive a private value, feasible bids and estimated win probabilities, removing forecasting and strategic learning. No securities, credit or banking data are used, and transfer to institutional finance is untested.

The evaluation contribution combines held-out bidder profit and regret, a construction-calculator comparator, observable objective-agreement classes, output-validity and aggregation checks, and recorded inference-call resources. It connects numerical consistency to decision consequences without treating either fluent explanations or conditional-on-valid-output scores as sufficient workflow validation. The ancillary service contest is analyzed in Appendix~\ref{app:servicecontest}: it tests allocation assumptions using the measured frontiers, rather than evaluating interacting LLM agents or governance interventions.
\section{Related work}
Test-time scaling studies compare prompt, sampling and search budgets \citep{wei2022chain,snell2024scaling,wu2024inference,alomrani2025budget}. Centralized allocation \citep{wan2026shadowprice} and confidence-based stopping \citep{jurayj2025finalanswer} ask different questions from service competition. Economic benchmarks study rational choice, auctions and interacting markets \citep{fish2025econevals,chen2024aucarena,zhu2024synthetic}; Bazaar studies repeated pricing with hidden preferences \citep{ahmed2026bazaar}. The contest in Appendix~\ref{app:servicecontest} draws on rent seeking \citep{tullock1980efficient} and compute-provider games \citep{artola2026games}.
The main empirical addition is held-out economic scoring with a construction-calculator comparator, observable action-loss diagnostics and resource accounting. Appendix~\ref{app:servicecontest} uses that evidence in a conditional service-allocation sensitivity analysis; neither component evaluates an autonomous multi-agent market.
% TODO: add 2-4 references on LLM numerical/financial reasoning evaluation
\section{Economic decision measurement}
\label{sec:framework}
\subsection{Economic consequences of a bid choice}
A risk-neutral bidder with value $v$ selects a feasible $b\leq v$. Historical closing price $P$ is an exogenous price-to-beat proxy. Its empirical distribution gives
\begin{equation}\begin{aligned}
\widehat p(b)&=\widehat F_P(b)=N^{-1}\sum_{n=1}^N\mathbf1\{P_n\leq b\},\\
U(b\mid v)&=\widehat p(b)(v-b).
\end{aligned}\label{eq:auctionutility}\end{equation}
The first-price payoff is a constructed decision objective; historical closing prices do not identify a first-price strategic equilibrium. Let $b_s^*=\arg\max_{b\in B_s}U_s(b\mid v_s)$ and let $\widehat b_{msc}$ be model $m$'s choice at budget $c$:
\begin{equation}\begin{aligned}
A_{msc}&=\mathbf1\{\widehat b_{msc}=b_s^*\},\\
R_{msc}&=U_s(b_s^*\mid v_s)-U_s(\widehat b_{msc}\mid v_s),\\
Q_{msc}&=U_s(\widehat b_{msc}\mid v_s)/U_s(b_s^*\mid v_s).
\end{aligned}\label{eq:metrics}\end{equation}
All optimum profits are positive. Repetitions are averaged within scenario/budget cells. Equal-scenario summaries are
\begin{equation}\begin{aligned}
q_m(c)&=S^{-1}\sum_s Q_{msc},\qquad
\overline R_m(c)=S^{-1}\sum_s R_{msc},\\
\overline A_m(c)&=S^{-1}\sum_s A_{msc}.
\end{aligned}\label{eq:frontier}\end{equation}
Mean dollar profit $\overline U_m(c)=S^{-1}\sum_s U_s(\widehat b_{msc}\mid v_s)$ instead weights each dollar equally. It need not rank intermediate budgets like mean scenario ratios. Neither statistic is total auction surplus. We retain the discrete frontiers without smoothing.

\section{Data, inference protocol and uncertainty}
\label{sec:protocol}
We use 10,681 bid records from 628 eBay auctions in Palm Pilot M515 PDA, Xbox and Cartier categories \citep{jank2010modeling,onlineauctionsdataset}. After excluding inconsistent opening bids, seed 20260819 assigns auctions within product to a 70/30 construction/scoring split. The 216 constructed tasks form 54 product/private-value families, with four candidate variants across five- and seven-bid menus. Candidate generation does not condition on the optimum.

Prompts supply private value, bids and construction-sample win probabilities rounded to four decimals; scoring uses held-out prices. The calculator maximizes the construction objective; a display-matched sensitivity checks the effect of its precision advantage. Option-order replications retain labels rather than positions. This setup isolates calculation/action selection, not financial forecasting, strategic bidding or elicited preferences.

Six checkpoints span Phi-mini, Qwen3-4B/8B/14B, DeepSeek-8B and Granite-8B. Greedy assigned budgets are $\{0,\allowbreak64,\allowbreak128,\allowbreak256,\allowbreak512,\allowbreak768\}$, with orders $\{11,23,37\}$. The stochastic frontier protocol uses $\mathcal C_S=\{128,\allowbreak384,\allowbreak768\}$, five seeds $\{1,7,21,42,99\}$, temperature 0.6, top-$p$ 0.95 and fixed order 314159. For positive budgets, the runner fixes minimum and maximum new tokens and suppresses EOS. Thus budget is a token boundary, not semantic completion. A separate constrained-label forward pass selects the candidate; final-answer/readout tokens are outside $c$. No-thinking, native-length and free-form controls are separate. Granite has a different native template. The original four checkpoint revisions were attributed retrospectively from a disclosed server-cache audit; Qwen3-14B and Granite revisions were recorded contemporaneously.

Greedy endpoints average orders before 2,000 paired scenario draws. The same sampled scenarios retain cross-budget dependence:
\begin{equation}
q_m^{(d)}(c)=S^{-1}\sum_{r=1}^S Q_{m,I_r^{(d)},c}.
\label{eq:bootstrapfrontier}\end{equation}
Economic contrasts use 10,000 product-stratified value-family draws, preserving all four variants. These exploratory intervals condition on the original products, prompts and construction/scoring split. A separate auction-CDF bootstrap re-scores fixed choices; it does not represent every source of uncertainty. Appendix~\ref{app:servicecontest} reports joint runtime/quality uncertainty for the cost-scenario sensitivity, including its ratio-conditioning rules.

\section{Evaluation design for numerical decision services}\label{sec:evaluation}
\subsection{Held-out utility and construction-calculator comparison}
All six greedy 0-to-768 utility gains are positive (Table~\ref{tab:endpoints}, Figure~\ref{fig:benchmark}); intermediate curves have dips and plateaus. Five stochastic generations retain positive 128-to-768 endpoint gains of 0.067--0.189. Family-clustered dollar-profit gains are \$11.11--\$46.36 per constructed decision, with pointwise intervals above zero. Greedy and stochastic estimates are separate protocols, not interchangeable endpoints.
\input{tables/endpoints}
\begin{figure*}[t]\centering
\includegraphics[width=0.9\textwidth]{figures/benchmark_frontiers.pdf}
\Description{Left: stochastic utility curves at three assigned budgets for six checkpoints. Right: greedy zero-to-768 aggregate calculator-gap closure with conditional family intervals.}
\caption{Separate protocols: left shows means over 216 scenarios and five stochastic seeds at 128, 384 and 768 tokens; right shows order-averaged greedy gap closure, with pointwise 95\% value-family bootstrap intervals. Gap closure is a fraction of aggregate mean-ratio deficit, not the fraction of tasks solved. The calculator remains stronger in simultaneous dollar-profit comparisons.}\label{fig:benchmark}\end{figure*}

The full-precision construction calculator has $q_{\rm calc}=0.981364$. The gap-closed statistic is $(q_m(768)-q_m(0))/(q_{\rm calc}-q_m(0))$, an aggregate mean-ratio comparison, not the fraction of tasks solved. At 768 tokens, all six model-minus-calculator dollar-profit simultaneous bands remain negative. Rounding the calculator's probabilities to the prompt's display precision changes one scenario's bid and mean profit by only \$0.0113, retaining the six disadvantages. More reasoning therefore shows progress without eliminating the comparator deficit.

The full-precision construction calculator, scored on held-out prices, earns mean expected bidder profit of \$212.66 and mean held-out regret of \$3.09; it is utility-optimal on 155 of 216 scoring scenarios. Table~\ref{tab:calculatorerroraudit} reports the greedy 768-token comparison, with exploratory simultaneous 95\% bands across six checkpoints. Each row retains all 216 scenarios and three option orders. The fraction of records beating the construction calculator on held-out profit is only 2.2--8.0\%; occasional wins do not reverse the aggregate deficit.
\input{tables/calculator_error_audit}
\subsection{Three-part loss and resource audit}
A numerical decision service should be evaluated on three distinct components, rather than a single success score:
\begin{enumerate}
\item \textbf{Action loss under displayed information.} Let $b_{\rm disp}^*$ maximize expected profit using the four-decimal probabilities actually shown to the model. Displayed-objective regret compares the chosen bid with that optimum under the same information. At 768 tokens, 42.1--86.3\% of choices maximize the displayed objective; 11.9--52.0\% are suboptimal under both displayed and held-out objectives. These observable action classes audit numerical consistency, not the model's internal arithmetic or reasoning mechanism.
\item \textbf{Distribution-transfer loss.} Displayed-optimal but held-out-suboptimal actions constitute 9.4--22.8\% of 768-token choices. They illustrate why optimizing the supplied construction distribution need not optimize held-out profit. This class is a diagnostic of objective agreement, not a causal attribution to distribution shift alone: rounding and ties also enter the comparison.
\item \textbf{Complete recorded service-call resources.} Inference-call accounting should cover generation plus answer, constrained readout, and confidence queries where the protocol performs them. The historical counters record these stages separately. Their sums exclude loading, failed/preliminary work, downtime, some host overhead and project-wide wall clock; they are not complete deployment cost or matched GPU/CPU latency. Section~\ref{sec:callresources} states the measured boundary.
\end{enumerate}
These are three audit components, not three independent, nonnegative terms that can simply be added. The held-out regret accounting is
\[
R^{\rm test}(\widehat b)=R^{\rm test}(b_{\rm disp}^*)+
\bigl[U^{\rm test}(b_{\rm disp}^*)-U^{\rm test}(\widehat b)\bigr].
\]
The second term can be negative. Displayed-objective loss uses a different payoff function; its magnitude is not interchangeable with that held-out difference. Resource use is reported alongside decision loss rather than silently converted into bidder profit. This separation makes the evaluation design auditable while preserving the limits of what the experiment identifies.

\begin{table*}[t]
\caption{Reporting checklist for numerical decision services.}\label{tab:reportingchecklist}
\centering\small
\begin{tabularx}{\textwidth}{@{}p{0.20\textwidth}XX@{}}\toprule
Audit item & What we measured here & What remains unmeasured \\\midrule
Deterministic comparator & Construction calculator and display-matched sensitivity versus model choices & Tool routing and escalation policies \\
Displayed-objective optimality & Agreement with the optimum under supplied, displayed probabilities & Internal arithmetic and reasoning mechanisms \\
Held-out transfer loss & Objective disagreement and regret on held-out auction prices & Causal attribution to distribution shift alone; financial transfer \\
Ratio vs dollar aggregation & Scenario ratios, dollar profit, concentration and corrected adjacent contrasts & Broad-population or causal decomposition \\
Offered-decision denominator (validity) & Free-form validity, paired-valid utility sensitivities and complete constrained-label decisions & Utility under an invalid-action, abstention or retry policy \\
Recorded stage resources & Generation/answer, readout and confidence-query counters, with protocol-specific coverage & Loading, failed work, complete wall clock and end-to-end deployment cost \\
\bottomrule\end{tabularx}
\end{table*}

\subsection{Ratio versus dollar aggregation}
Nine of thirty adjacent contrasts (five adjacent budget pairs across six checkpoints) have opposite point directions under ratio and dollar aggregation. Zero of those nine have both directions significant after Holm correction across sixty exploratory family sign-flip tests; those tests assume sign symmetry under the null. The largest 22 tasks contribute 80.0--86.7\% of positive dollar gains; negative gains remain in the means. Product, option width and checkpoint main effects explain little of this variation. These are descriptive concentration checks, not a causal attribution or broad-population decomposition.

\subsection{Invalid outputs and decision denominators}
The fixed free-form parser accepts displayed candidate bids only; invalid responses remain missing, and validity ranges from 29.1\% to 96.7\%. Paired-valid sensitivities change sample composition. Readout disagreement is associated with greater regret for five checkpoints at 768; Granite's contrast has only twelve disagreement records and is excluded from pooled statements. Literal closure markers cannot establish semantic completion or the causal effect of forced length.
For workflow-level evaluation, report validity and utility on an \emph{offered-decision denominator}, alongside any conditional-on-valid-output statistics. This avoids presenting a selected subset as the performance of all offered decisions. A prospective service must specify how invalid actions are handled; an abstention payoff or retry policy is not invented here. The free-form utility sensitivities remain conditional on valid paired cells. The primary constrained-label results remain separate and retain all candidate decisions; this recommendation does not rescore either protocol.

\subsection{Construction calculator as a required tool comparator}
The construction calculator is a mandatory comparator for this benchmark because the supplied quantities define a small numerical optimization. The endpoint gains establish progress with assigned reasoning budget, not a reason to replace that comparator: it remains stronger in simultaneous dollar-profit comparisons, and intermediate gains are uneven. This is an evaluation of a comparator, not an experiment in which an LLM calls a tool.

\textbf{Untested design implication.} A potential decision-support workflow should retain a deterministic calculation path when the objective and numerical inputs are known, and evaluate the LLM on language interpretation or uncertain-context components. That division of labor, tool routing and any escalation policy have not been implemented or evaluated. Before applying the design to investment, credit or risk, one would need an appropriate action set, payoff, information-at-decision-time specification, financial data and temporal validation. Table~\ref{tab:reportingchecklist} summarizes what this benchmark does and does not measure.

\subsection{Recorded service-call resources}\label{sec:callresources}
The audit retains 50,544 historical records, 8,424 per checkpoint, and separates primary generation/answer, constrained readout and verbal-confidence query counters (Table~\ref{tab:resources}). All recorded primary timers are finite. Loading, downtime, host overhead, failed/preliminary work and complete project wall clock remain outside these sums. Whole-device energy telemetry is not isolated-model energy. These stage totals describe the historical protocols, not matched-hardware speed rankings.
\input{tables/resources}
For the fixed-budget protocol, the primary timer already includes generation and answer, so those stages must not be counted twice when adding constrained-readout and confidence-query counters. Stage coverage varies across historical control protocols and must be checked against the runner before summing components. Table~\ref{tab:resources} preserves the primary/readout totals; the separately stored confidence-query counters are retained in the companion evidence. Recorded stage sums are a resource-accounting component of the evaluation, not an end-to-end service-cost estimate. A complete deployed-call measurement would also need the excluded overheads and its actual invalid-output/retry policy.

\section{Resource-cost scenarios and incentive sensitivity (summary)}\label{sec:costsummary}
Hardware and prices are cost scenarios, not historical bills or a controlled comparison across all checkpoints. An assumed service contest, rather than observed LLM-agent behavior, maps quality into private and planner allocations. At $\gamma=1$, pooled private payment break-even values are \$0.03--\$0.18 and planner coefficients are \$0.02--\$0.12 per unit benchmark quality. Qwen3-14B uses a separate L40S scenario; the other five use RTX 4090 scenarios. Product-specific private bounds vary up to 30.40-fold, so the pooled headline is task-mixture dependent. These bounds support only the maximum \emph{tested} budgets, not an unrestricted optimum.

Appendix~\ref{app:servicecontest} gives the model, complete finite-game analysis, proofs, cost accounting and the break-even table (Table~\ref{tab:money}). It is an incentive sensitivity analysis, not evidence of strategic compute choices or identified client demand.

\section{Limitations and reproducibility}\label{sec:limitations}
Three products, fixed prompts/splits, supplied probabilities and a small bid menu limit transfer. Historical prices are an exogenous proxy. Forced tokens imperfectly represent useful computation, templates differ, and original-four revision attribution is retrospective. A development diagnostic/probe uses only six adapted families and supports no primary claim; its richer raw-record extraction losses do not substitute for independent validation. Human-study materials have no collected responses.

Independent richer/strategic decisions, matched native and above-768 reasoning, observed client selection/value and endogenous compute choices remain missing. Modeled service welfare (Appendix~\ref{app:servicecontest}) excludes latency/environmental externalities and heterogeneous clients. Saved-output replay cannot fill those identification gaps. The anonymous primary release preserves prompts, revision provenance and full historical outputs; the companion evidence package preserves CPU scripts, compact inputs, hash manifests and conditional uncertainty.\footnote{\url{https://anonymous.4open.science/r/reasoning-compute-economic-utility-EED7/}} All analyses reuse the saved runs; no new GPU inference was performed. The main evaluation results and limits are stated here; the ancillary contest assumptions and results are reported in Appendix~\ref{app:servicecontest}. The companion evidence supports saved-data replay without adding independent confirmation.
\section{Conclusion}
Held-out economic utility, a construction-calculator comparator and recorded resources provide complementary evidence for evaluating numerical LLM decision services. Assigned reasoning budgets improve endpoint utility while leaving an aggregate calculator deficit; weighting, invalid-output denominators and task composition materially affect interpretation. The design implications for deterministic calculation paths and financial workflows remain untested. Appendix incentive sensitivity clarifies compute assumptions without establishing observed strategic behavior or financial deployment.
\FloatBarrier
\bibliographystyle{ACM-Reference-Format}
\bibliography{references}
\clearpage
\appendix
\section{Service-contest sensitivity model}\label{app:servicecontest}
This appendix presents the ancillary allocation model and its sensitivity evidence. It is an assumed contest between recommendation providers, not observed LLM-agent behavior or identified financial service demand.
\input{tables/money}
\subsection{A specified contest between recommendation providers}
Two providers choose $c_i\in\mathcal C$ and pay incremental inference cost. Benchmark quality is treated as a proxy for recommendation quality; it is not empirically identified service demand. The selected provider receives fixed payment $V$. For symmetric quality $q(c)>0$, the Tullock rule and private payoff are
\begin{equation}\begin{aligned}
P_i(c_i,c_j;\gamma)&=\frac{q(c_i)^\gamma}{q(c_i)^\gamma+q(c_j)^\gamma},\\
\pi_i(c_i,c_j)&=VP_i(c_i,c_j;\gamma)-\lambda c_i.
\end{aligned}\label{eq:contest}\end{equation}
Here $\gamma>0$ is specified selection sensitivity; service selection differs from the bid-win probability $\widehat p$. Dividing by payment gives
\begin{equation}
\widetilde\pi_i=P_i-\kappa c_i,\qquad \kappa=\lambda/V.
\label{eq:normalizedprivate}\end{equation}
If selected-service benefit is $\beta q(c_{\rm winner})$, contract payments cancel when client and provider payoffs are summed:
\begin{equation}\begin{aligned}
W_\beta(c_1,c_2)&=\beta Q_\gamma(c_1,c_2)-\lambda(c_1+c_2),\\
Q_\gamma(c_1,c_2)&=P_1q(c_1)+P_2q(c_2).
\end{aligned}\label{eq:welfare}\end{equation}
This modeled service welfare differs from auction surplus. Its normalized form is
\begin{equation}
\widetilde W=\rho Q_\gamma-\kappa(c_1+c_2),\qquad \rho=\beta/V.
\label{eq:normalizedwelfare}\end{equation}
A symmetric pure equilibrium is a mutual best response; we select the cheapest symmetric one under ties and separately enumerate the full equilibrium set. The planner maximizes over all ordered pairs. Total private compute above, equal to or below planner compute defines overcompute, alignment or undercompute. Against a fixed rival action, an increase from $c_a$ to $c_b$ is privately attractive when
\begin{equation}
\kappa<\frac{P_i(c_b,c_j;\gamma)-P_i(c_a,c_j;\gamma)}{c_b-c_a}.
\label{eq:privatethreshold}\end{equation}
For planner changes $(\Delta Q_\gamma,\Delta C)$, the comparison is
\begin{equation}
\kappa/\rho=\lambda/\beta<\Delta Q_\gamma/\Delta C.
\label{eq:plannerthreshold}\end{equation}
For provider-specific frontiers and costs, substitute $q_i(c_i)$ and $\lambda_i$; cross-checkpoint quality comparability remains an assumption. Absolute modeled welfare loss is $\Delta W=W_{\rm OPT}-W_{\rm NE}$. Relative loss divides by $W_{\rm OPT}$ only when positive; otherwise it is undefined, without a near-zero cutoff.
\subsection{Finite-game enumeration}
All nine planner pairs and all private pure best responses are enumerated. For the symmetric game, let $R_{ab}=P_{ab}-1/2-\kappa(c_a-c_b)$. Adding the rival's cost leaves best responses unchanged and gives a constant-sum game. $R$ is skew symmetric, with value zero and common optimal-strategy polytope
\begin{equation}
\mathcal S=\{x\geq0:\mathbf1^\top x=1,\ Rx\leq0\}.
\label{eq:equilibriumpolytope}\end{equation}
The complete Nash set is $\mathcal S\times\mathcal S$; vertices and vertex pairs bound expected compute and welfare. Exact payoff crossings define allocation boundaries; phase-grid areas have no market-frequency interpretation. Ties use cheapest best responses and lexicographic planner maximizers.
\subsection{Cost-accounting scenarios}
Incremental costs use OLS runtime slopes with an intercept over the three stochastic budgets. The five RTX 4090 scenarios use \$0.74/hour; Qwen3-14B uses its own L40S slope at a separate \$1/hour scenario \citep{runpod2026pricing}. These imply \$2.46--\$4.46$\times10^{-6}$ and \$1.23$\times10^{-5}$ per reasoning token, respectively. Hardware is not matched across all six models, and the prices are accounting scenarios rather than measured historical bills.

An electricity-only component uses incremental whole-device joules at hypothetical \$0.10/kWh. It is neither a complete deployment-cost bound nor added to rental accounting. Illustrative \$10$^{-5}$ and \$10$^{-4}$/token brackets are not measured API tariffs. At the higher bracket, pooled private-support margins at secondary $V_0=\$328.23$ remain 110--284. A local 21,600-call calculator-kernel benchmark has median 1.917 microseconds on a different CPU host, excludes parsing, and supports no GPU/CPU cost or speed comparison.
\subsection{Joint runtime and quality uncertainty}
Economic-contrast and auction-CDF bootstrap details are reported above (see Sec.~\ref{sec:protocol}). Joint cost/quality draws (seed 20261011) sample a stored seed for each selected scenario occurrence jointly across budgets, quality, runtime and energy. Prices and hardware are held fixed; these intervals do not capture deployment uncertainty. These are exploratory conditional intervals, fixing products, prompts, construction/scoring split, hardware and price. Ratio intervals condition on positive support bounds: Phi has seven of 10,000 draws with nonpositive private and planner bounds, retained as unsupported counts. Absolute-bound draws remain included.

\subsection{Exact boundaries and break-even values}\label{sec:allocation}
\begin{proposition}[Maximum-budget support]\label{prop:maxsupport}
Let $c_H$ be the maximum budget. For symmetric positive quality under the stated selection rule, define
\[
K=\min_{c<c_H}\frac{P(c_H,c_H)-P(c,c_H)}{c_H-c}.\]
\[
r_1=\min_{(a,b)\ne(H,H)}\frac{q(c_H)-Q_\gamma(c_a,c_b)}{2c_H-c_a-c_b}.
\]
If $K,r_1>0$, high/high is a pure equilibrium iff $V\geq\lambda/K$, and is a planner maximizer iff $\beta\geq\lambda/r_1$, allowing equality ties.
\end{proposition}
\begin{proof}
For each unilateral lower-budget deviation, high/high support is $\lambda/V\leq[P(H,H)-P(c,H)]/(c_H-c)$. Taking the minimum gives $K$. Comparing high/high with every lower-total planner pair gives $\lambda/\beta\leq[q(c_H)-Q_\gamma(a,b)]/(2c_H-c_a-c_b)$, hence $r_1$. Rearranging proves both bounds. If a bound is nonpositive, that ratio cannot give positive-cost support.
\end{proof}
At $\gamma=1$, pooled private break-evens $V^*=\lambda/K$ are \$0.034--\$0.080 for the five RTX scenarios and \$0.176 for Qwen3-14B's own L40S scenario. Planner coefficients $\beta^*=\lambda/r_1$ are \$0.022--\$0.051 and \$0.116 per unit benchmark quality (Table~\ref{tab:money}). These avoid dependence on an arbitrary reference payment. Above the respective bounds, both choices retain the maximum tested budgets; alignment is constrained by the menu, not proof of an unrestricted optimum.

At $\gamma=1$, every pooled stochastic frontier yields overcompute, alignment and undercompute regions. Four checkpoints jump from 768 to 128 as private normalized cost rises; Qwen3-4B and Granite pass through 384. Planner total compute follows $1536\to896\to256$; the middle pair concentrates 768 on one provider and 128 on the other.

\begin{proposition}[Structural regions]\label{prop:regions}
For a finite menu with distinct costs, positive quality strictly maximized at its highest budget, and finite $\gamma>0$, all three allocation regions occur on open neighborhoods when positive cost and client-value domains are sufficiently broad.
\end{proposition}
\begin{proof}
At zero cost, highest quality strictly dominates under the Tullock rule, and remains dominant at sufficiently small positive cost. Bounded selection gains make the cheapest action dominant at sufficiently large private cost. At fixed positive cost, the planner selects the cheapest pair for sufficiently small client value and highest/highest for sufficiently large value. Combining extremes yields all three regions; finiteness and strict choices give open neighborhoods. A flat frontier instead yields cheapest-action alignment at positive cost.
\end{proof}
The regions' existence is structural; the data locate their boundaries. Across six frontiers and four tested $\gamma$ values, full finite-game analysis finds one mixed-only open interval: Qwen3-4B at $\gamma=4$, approximately $(3.28603,3.28909)\times10^{-4}$. Score-origin shifts, isotonic envelopes and matched logit selection expose mapping dependence. Accuracy-based private support thresholds are 4.05--5.96 times utility-based bounds, with all six simultaneous differences positive. These changes concern selection assumptions, not measured client choice. Product-specific private bounds vary 1.59--30.40-fold within checkpoint, with Granite at the upper end: pooled cents-level thresholds depend on task mixture.

\subsection{Different providers and modeled welfare}
Fifteen distinct pairs at a common cost and ten RTX pairs at their own costs are evaluated at seven payments and $\gamma\in\{1,4\}$ with $\beta=V$: 210 common-cost plus 140 own-RTX games. Each of these 350 sampled points has one pure equilibrium, without establishing uniqueness over all parameters or excluding mixed equilibria. For DeepSeek/Phi at own costs and $\gamma=1$, payment \$0.03 gives private 128/128 and planner 768/128; \$0.07 gives private 768/128 and planner 768/768. Welfare losses are 2.69\% and 0.50\% of positive planner welfare, respectively. All 350 planner welfare denominators are positive. Changing $V$ alone leaves the planner unchanged at fixed $\beta$; a path with $\beta=V$ changes client value too.
\end{document}
